\documentclass{article}
\usepackage[utf8]{inputenc}
\usepackage{amsmath}
\usepackage{lmodern,amsthm,amsfonts,amssymb,anyfontsize,thm-restate}

\newtheorem{innercustomthm}{Theorem}
\newtheorem{innercustomlem}{Lemma}
\newtheorem{definition}{Definition}
\newtheorem{innercustomdef}{Definition}
\newtheorem*{definition*}{Definition}
\newtheorem*{thm*}{Theorem}
\newtheorem{exercise}{Exercise}[section]
\newtheorem*{claim}{Claim}
\title{A whole lot of reductions}
\author{Sergio J. Hernando Rivero}
\date{\today}

\begin{document}

\maketitle

\section{preliminaries}
\begin{definition}
    A \textbf{tour} in a graph $G = (V, E)$ is a sequence of vertices and edges adjacent to them in the graph that spans every vertex. We usually refer to the multiset of edges defined by a tour as $T$. The \textbf{cost} of a tour $T$ is the sum of the weights of the edges in $T$, denoted as $c(T) = \sum_{e \in T} c(e)$.
\end{definition}
This definition is equivalent to saying that a tour is a Eulerian multigraph that spans every vertex of $G$.
\begin{definition}[Metric TSP]
    Let $G = (V, E)$ be a connected undirected graph. The \textbf{Metric Traveling Salesperson Problem} (Metric TSP) is defined as follows:
    \begin{itemize}
        \item \textbf{Instance:} A complete graph $G = (V, E)$ which is connected by definition, and a cost function $c: E \to \mathbb{R}_{\geq 0}$ that satisfies the triangle inequality:
        \begin{equation}
            c(u, v) \leq c(u, w) + c(w, v) \quad \text{for all } u, v, w \in V.
        \end{equation}
        \item \textbf{Objective:} Find a closed tour $C$ in $G$ such that the total weight $c(C) = \sum_{e \in C} c(e)$ is minimized.
    \end{itemize}
\end{definition}
\begin{definition}[Quasi-Tour]
    A \textbf{quasi-tour} $T_0$ in a graph $G = (V, E)$ is a subset of edges in $T$ removing edges from a finite set of closed tours. 
    Each connected component of this multigraph is either a Eulerian multigraph or an isolated vertex. 
\end{definition}
A given multiset of edges $T_0$ is a quasi-tour if and only if it is Eulerian.

\section{Starting point}
\begin{definition}[MAX-E3-LIN-2]
    The \textbf{MAX-E3-LIN-2} problem is defined as follows:
    \begin{itemize}
        \item \textbf{Instance:} A set of $n$ boolean variables $x_1, x_2, \ldots, x_n$ and a collection of $m$ linear equations over $\mathbb{F}_2$, where each equation involves exactly three variables. Each equation is of the form:
        \begin{equation}
            x_{i_1} + x_{i_2} + x_{i_3} = b_j,
        \end{equation}
        where $b_j \in \{0, 1\}$ and $i_1, i_2, i_3$ are distinct indices from $\{1, 2, \ldots, n\}$.
        \item \textbf{Objective:} Find an assignment to the variables that maximizes the number of satisfied equations.
    \end{itemize}
\end{definition}
\end{document}